\chapter{The Riemann Zeta Function}
  Throughout, \(s = \s+it\) will stand for complex variables with \(\s\) and \(t\) real. We will also make use of the gamma function developed in \cref{sec:the_gamma_function}. Also, all sums and products over zeros of a function will be counted with multiplicity and ordered symmetrically with respect to the size of the ordinate if they are not absolutely convergent.
  \section{\todo{Analytic Properties}}
    The \textbf{Riemann zeta function}\index{Riemann zeta function} \(\z(s)\) is defined by the Dirichlet series
    \[
      \z(s) = \sum_{n \ge 1}\frac{1}{n^{s}}.
    \]
    As the coefficients are uniformly bounded and completely multiplicative, \(\z(s)\) is absolutely convergent for \(\s > 1\) and admits the following degree \(1\) Euler product
    \[
      \z(s) = \prod_{p}(1-p^{-s})^{-1},
    \]
    by \cref{prop:Dirichlet_series_Euler_product}.

    Our primary aim is to show that the Riemann zeta function is a Selberg class \(L\)-function. In order to do this, we need to establish meromorphic continuation to the entire complex plane and prove a functional equation. These will be accomplished in tandem by developing and analyzing an integral representation for the Riemann zeta function that evolves from a Mellin transform. The function which we will take the Mellin transform of is closely related to the \textbf{Jacobi theta function}\index{Jacobi theta function} \(\vt(z)\) defined by
    \[
      \vt(z) = \sum_{n \in \Z}e^{2\pi in^{2}z},
    \]
    for \(z \in \H\). This function is locally absolutely uniformly convergent in this region by the Weierstrass \(M\)-test. Moreover, the sum of a geometric series implies
    \[
      \vt(z)-1 = O(e^{-2\pi y}),
    \]
    so that \(\vt(z)-1\) exhibits exponential decay. The essential result we will need about the Jacobi theta function is that it satisfies a transformation property. This will follow by an application of Poisson summation.

    \begin{theorem}\label{thm:functional_equation_Jacobi_theta}
      For \(z \in \H\),
      \[
        \vt(z) = \frac{1}{\sqrt{-2iz}}\vt\left(-\frac{1}{4z}\right).
      \]
    \end{theorem}
    \begin{proof}
      We will apply Poisson summation to the Jacobi theta function. By the identity theorem it suffices to verify the transformation formula for \(z = iy\) with \(y\) positive. So set
      \[
        f(x) = e^{-2\pi x^{2}y}.
      \]
      Then \(f(x)\) is Schwartz class. By \cref{prop:Fourier_transform_of_exponential_single_variable}, we have
      \[
        (\mc{F}f)(t) = \frac{e^{-\frac{\pi t^{2}}{2y}}}{\sqrt{2y}}.
      \]
      Poisson summation together with this identity gives the result.
    \end{proof}

    We can now introduce the function whose Mellin transform will essentially give the Riemann zeta function. This function \(\w(z)\) is defined by
    \[
      \w(z) = \sum_{n \ge 1}e^{\pi in^{2}z},
    \]
    for \(z \in \H\). It is related to the Jacobi theta function via
    \[
      \w(z) = \frac{\vt\left(\frac{z}{2}\right)-1}{2},
    \]
    and hence is locally absolutely uniformly convergent in this region. In particular, we have the estimates
    \[
      \w(iy) = \begin{cases} O\left(y^{-\frac{1}{2}}\right) & \text{if \(y \ll 1\)}, \\ O(e^{-\pi y}) & \text{if \(y \gg 1\)}, \end{cases}
    \]
    where the second estimate follows from the exponential decay of the Jacobi theta function and the first estimate from \cref{thm:functional_equation_Jacobi_theta}. Now consider the Mellin transform
    \[
      \int_{0}^{\infty}\w(iy)y^{\frac{s}{2}}\,\frac{dy}{y}.
    \]
    By the previous estimates, this integral is locally absolutely uniformly convergent for \(\s > 1\) and hence defines an analytic function there. This also permits the us to interchange the integral and the sum inside the integrand resulting in
    \[
      \int_{0}^{\infty}\w(iy)y^{\frac{s}{2}}\,\frac{dy}{y} = \sum_{n \ge 1}\int_{0}^{\infty}e^{-\pi n^{2}y}y^{\frac{s}{2}}\,\frac{dy}{y}.
    \]
    Under the change of variables \(y \mapsto \frac{y}{\pi n^{2}}\), the integral is realized as \(\G\left(\frac{s}{2}\right)\) and the sum becomes \(\z(s)\) giving
    \[
      \sum_{n \ge 1}\int_{0}^{\infty}e^{-\pi n^{2}y}y^{\frac{s}{2}}\,\frac{dy}{y} = \frac{\G\left(\frac{s}{2}\right)}{\pi^{\frac{s}{2}}}\z(s).
    \]
    As the reciprocal of the gamma function is entire, we obtain an integral representation
    \begin{equation}\label{equ:integral_representation_zeta_1}
      \z(s) = \frac{\pi^{\frac{s}{2}}}{\G\left(\frac{s}{2}\right)}\int_{0}^{\infty}\w(iy)y^{\frac{s}{2}}\,\frac{dy}{y}.
    \end{equation}
    This integral representation does not establish the meromorphic continuation of the Riemann zeta function to the entire complex plane because the integral does not converge for \(\s \le 1\) in the region where \(y \ll 1\) in view of the estimates we have just establishes. In order to bypass this obstruction, we must decompose the integral.
    \iffalse
    Returning to the Riemann zeta function, we split the integral in \cref{equ:integral_representation_zeta_1} into two pieces by writing
    \begin{equation}\label{equ:symmetric_integral_zeta_split}
      \int_{0}^{\infty}\w(iy)y^{\frac{s}{2}}\,\frac{dy}{y} = \int_{0}^{1}\w(iy)y^{\frac{s}{2}}\,\frac{dy}{y}+\int_{1}^{\infty}\w(iy)y^{\frac{s}{2}}\,\frac{dy}{y}.
    \end{equation}
    The idea now is to rewrite the first piece in the same form and symmetrize the result as much as possible. We being by performing a change of variables \(y \mapsto \frac{1}{y}\) to the first piece to obtain
    \[
      \int_{1}^{\infty}\w\left(\frac{i}{y}\right)y^{-\frac{s}{2}}\,\frac{dy}{y}
    \]
    Now we compute
    \begin{align*}
      \w\left(\frac{i}{y}\right) &= \w\left(-\frac{1}{iy}\right) \\
      &= \frac{\vt\left(-\frac{1}{2iy}\right)-1}{2} \\
      &= \frac{\sqrt{y}\vt\left(\frac{iy}{2}\right)-1}{2} && \text{\cref{thm:functional_equation_Jacobi_theta}} \\
      &= \sqrt{y}\w(iy)+\frac{\sqrt{y}}{2}-\frac{1}{2}.
    \end{align*}
    This chain implies that our first piece can be expressed as
    \[
      \int_{1}^{\infty}\left(\sqrt{y}\w(iy)+\frac{\sqrt{y}}{2}-\frac{1}{2}\right)y^{-\frac{s}{2}}\,\frac{dy}{y},
    \]
    which is further equivalent to
    \[
      \int_{1}^{\infty}\w(iy)y^{\frac{1-s}{2}}\,\frac{dy}{y}-\frac{1}{s(1-s)},
    \]
    because the integral over the last two pieces is \(\frac{1}{1-s}-\frac{1}{s} = -\frac{1}{s(1-s)}\). Substituting this result back into \cref{equ:symmetric_integral_zeta_split} and combining with \cref{equ:integral_representation_zeta_1} yields the integral representation
    \begin{equation}\label{equ:integral_representation_zeta_final}
      \z(s) = \frac{\pi^{\frac{s}{2}}}{\G\left(\frac{s}{2}\right)}\left[-\frac{1}{s(1-s)}+\int_{1}^{\infty}\w(iy)y^{\frac{1-s}{2}}\,\frac{dy}{y}+\int_{1}^{\infty}\w(iy)y^{\frac{s}{2}}\,\frac{dy}{y}\right].
    \end{equation}
    This integral representation will give analytic continuation. To see this, first observe that everything outside the brackets is entire. Moreover, the integrands exhibit exponential decay and therefore the integrals are locally absolutely uniformly convergent on \(\C\). The fractional term is holomorphic except for simple poles at \(s = 0\) and \(s = 1\). The meromorphic continuation to \(\C\) follows with possible simple poles at \(s = 0\) and \(s = 1\). There is no pole at \(s = 0\). Indeed, \(\G\left(\frac{s}{2}\right)\) has a simple pole at \(s = 0\) and so its reciprocal has a simple zero. This cancels the corresponding simple pole of \(-\frac{1}{s(1-s)}\) so that \(\z(s)\) has a removable singularity and thus is holomorphic at \(s = 0\). At \(s = 1\), \(\G\left(\frac{s}{2}\right)\) is nonzero and so \(\z(s)\) has a simple pole. Therefore \(\z(s)\) has meromorphic continuation to all of \(\C\) with a simple pole at \(s = 1\).
  \subsection*{The Functional Equation}
    An immediate consequence of applying the symmetry \(s \mapsto 1-s\) to \cref{equ:integral_representation_zeta_final} is the following functional equation:
    \[
      \frac{\G\left(\frac{s}{2}\right)}{\pi^{\frac{s}{2}}}\z(s) = \frac{\G\left(\frac{1-s}{2}\right)}{\pi^{\frac{1-s}{2}}}\z(1-s).
    \]
    We identify the gamma factor as
    \[
      \g(s,\z) = \pi^{-\frac{s}{2}}\G\left(\frac{s}{2}\right),
    \]
    with \(\k = 0\) the only local root at infinity. Clearly it satisfies the required bounds. The conductor is \(q(\z) = 1\) so no primes ramify. The completed Riemann zeta function is
    \[
      \L(s,\z) = \pi^{-\frac{s}{2}}\G\left(\frac{s}{2}\right)\z(s),
    \]
    with functional equation
    \[
      \L(s,\z) = \L(1-s,\z).
    \]
    This is the functional equation of \(\z(s)\) and in this case is just a reformulation of the previous functional equation. From it we find that the root number is \(\e(\z) = 1\) and that \(\z(s)\) is self-dual. We can now show that the order of \(\z(s)\) is \(1\). Since there is only a simple pole at \(s = 1\), multiply by \((s-1)\) to clear the polar divisor. As the integrals in \cref{equ:integral_representation_zeta_final} are locally absolutely uniformly convergent, computing the order amounts to estimating the gamma factor. Since the reciprocal of the gamma function is of order \(1\), we have
    \[
      \frac{1}{\g(s,\z)} \ll_{\e} e^{|s|^{1+\e}}.
    \]
    Thus the reciprocal of the gamma factor is also of order \(1\). It follows that
    \[
      (s-1)\z(s) \ll_{\e} e^{|s|^{1+\e}}.
    \]
    This shows \((s-1)\z(s)\) is of order \(1\) and thus \(\z(s)\) is as well after removing the polar divisor. We now compute the residue of \(\z(s)\) at \(s = 1\). First observe that the only term in \cref{equ:integral_representation_zeta_final} contributing to the pole is \(-\frac{\pi^{\frac{s}{2}}}{\G\left(\frac{s}{2}\right)}\frac{1}{s(1-s)}\). Then
    \[
      \Res_{s = 1}\z(s) = \Res_{s = 1}\left(-\frac{\pi^{\frac{s}{2}}}{\G\left(\frac{s}{2}\right)}\frac{1}{s(1-s)}\right) = \lim_{s \to 1}\left(\frac{\pi^{\frac{s}{2}}}{\G\left(\frac{s}{2}\right)}\frac{1}{s}\right) = 1,
    \]
    where the last equality follows from \cref{equ:gamma_function_at_1/2}. We summarize all of our work into the following theorem:

    \begin{theorem}\label{thm:zeta_Selberg}
      \(\z(s)\) is a Selberg class \(L\)-function with degree \(1\) Euler product given by
      \[
        \z(s) = \prod_{p}(1-p^{-s})^{-1}.
      \]
      Moreover, it admits meromorphic continuation to \(\C\), possesses the functional equation
      \[
        \pi^{-\frac{s}{2}}\G\left(\frac{s}{2}\right)\z(s) = \L(s,\z) = \L(1-s,\z),
      \]
      and has a simple pole at \(s = 1\) of residue \(1\).
    \end{theorem}

    Lastly, we note that by virtue of the functional equation we can also compute \(\z(0)\). Indeed, since \(\Res_{s = 1}\z(s) = 1\), we have
    \[
      \lim_{s \to 1}(s-1)\L(s,\z) = \left(\Res_{s = 1}\z(s)\right)\left(\lim_{s \to 1}\pi^{-\frac{s}{2}}\G\left(\frac{s}{2}\right)\right) = 1.
    \]
    In other words, \(\L(s,\z)\) has a simple pole at \(s = 1\) of residue \(1\) too. It follows that \(\L(s,\z)\) also has a simple pole at \(s = 0\) of residue \(1\). Hence
    \[
      1 = \lim_{s \to 1}(s-1)\L(1-s,\z) = \left(\Res_{s = 1}\G\left(\frac{1-s}{2}\right)\right)\left(\lim_{s \to 1}\pi^{-\frac{1-s}{2}}\z(1-s)\right) = -2\z(0),
    \]
    because \(\Res_{s = 0}\G(s) = 1\). Therefore \(\z(0) = -\frac{1}{2}\).
  \subsection*{The Inverse Riemann Zeta Function}
    It turns out that the inverse of the Riemann zeta function is the Dirichlet series whose coefficients are given by the M\"obius function:

    \begin{proposition}\label{prop:Dirichlet_Mobius_is_zeta_inverse}
      For \(\s > 1\),
      \[
        \z(s)^{-1} = \sum_{n \ge 1}\frac{\mu(n)}{n^{s}} = \prod_{p}(1-p^{-s}).
      \]
    \end{proposition}
    \begin{proof}
      Recall from \cref{prop:Mobius_indicator} that
      \[
        \mathbf{1} \ast \mu = \d.
      \]
      From \cref{equ:Dirichlet_convolution_of_Dirichlet_series}, we obtain
      \[
        \z(s)\sum_{n \ge 1}\frac{\mu(n)}{n^{s}} = 1,
      \]
      for \(\s > 1\). This is equivalent to first equality. The second follows from the Euler product for \(\z(s)\).
    \end{proof}
    \fi
  \section{\todo{The Prime Number Theorem}}
  \section{\todo{A Second Moment Estimate}}